\documentclass[a4paper, 11pt]{article}

% --- Packages ---
\usepackage[utf8]{inputenc}
\usepackage{graphicx}
\usepackage[T1]{fontenc}
\usepackage[french]{babel}
\usepackage[margin=1.5cm, top=2cm, bottom=2cm]{geometry}
\usepackage{xcolor}
\usepackage{fancyhdr}
\usepackage{lmodern}
\usepackage{titlesec}
\usepackage{enumitem}

% --- Définition des couleurs ---
\definecolor{primaryBlue}{RGB}{20, 50, 90}

% --- Styles de document ---
\pagestyle{fancy}
\fancyhf{}
\fancyfoot[L]{\footnotesize \textcolor{gray}{Projet Wakala --- Rapport d'Usabilité (Benchmark Concurrentiel)}}
\fancyfoot[R]{\footnotesize \textcolor{gray}{Page \thepage}}
\renewcommand{\headrulewidth}{0pt}
\renewcommand{\footrulewidth}{0.4pt}

% --- Personnalisation des titres ---
\titleformat{\section}{\Large\bfseries\color{primaryBlue}}{}{0em}{}[\titlerule]
\titleformat{\subsection}{\large\bfseries\color{primaryBlue}}{}{0em}{}

\begin{document}

% --- En-tête ---
\begin{center}
    
    {\Huge \textcolor{primaryBlue}{\textbf{WAKALA}}}\\[0.2cm]
    {\large \bfseries \textcolor{primaryBlue}{Rapport d'Usabilité : Analyse de la Concurrence}}
\end{center}

\vspace{0.5cm}

% --- Introduction ---
\section*{Introduction}
Quand on veut acheter ou vendre une voiture d'occasion sur internet, on cherche avant tout un outil simple, rapide et rassurant. Afin de comprendre les difficultés réelles que rencontre un utilisateur normal au quotidien, nous avons testé la navigation sur les principaux sites d'annonces au Maroc (\textit{Avito.ma}, \textit{Moteur.ma}) et à l'étranger (\textit{Carvana}, \textit{CarMax}). Ce document présente les 6 plus grandes frustrations et blocages ressentis au cours de ce test. L'objectif est d'éviter tous ces pièges sur notre propre plateforme \textbf{Wakala} afin de proposer l'expérience la plus facile et accessible possible.

\vspace{0.3cm}

% --- Les 6 Frictions ont été intégrées directement dans la synthèse ci-dessous ---

% --- Synthèse ---
\section*{Synthèse : Ce qu'il faut retenir}
Aujourd'hui, il y a beaucoup de choix sur les sites d'annonces automobiles, mais les utiliser reste très stressant et compliqué. Nous avons repéré quatre grands problèmes :

D'abord, la surcharge de publicités et de textes rend la navigation fatigante et fait perdre du temps. Ensuite, les recherches sont trop complexes : il faut souvent connaître des termes mécaniques très précis, ce qui bloque rapidement les utilisateurs novices. De plus, le manque de confiance règne en maître. Comme n'importe qui peut publier n'importe quoi sans vérification, la peur de l'arnaque ou des faux profils est permanente. Enfin, le parcours d'achat s'arrête en plein milieu : une fois la voiture trouvée, le site laisse l'utilisateur totalement seul pour négocier, gérer le paiement et faire la paperasse. Tout ce stress repose alors uniquement sur les épaules de l'acheteur.



\end{document}








